\section{Queues}

\solution{Problem}{43}
The solution is Listing~\ref{lst:twostack} of Chapter~\ref{ch:queues}; the amortised
bound is Theorem~\ref{thm:twostack}, and Tracing~\ref{trace:12} follows one run.
\idea \cc{in} receives new items; \cc{out} delivers the oldest. Refill \cc{out} from
\cc{in} only when \cc{out} is empty: the transfer reverses the order once, which turns
LIFO into FIFO.
\complexity $\Theta(1)$ amortised per operation, $\Theta(n)$ for a single dequeue that
triggers a transfer.
\mistakes Transferring whenever \cc{in} is non-empty (older items in \cc{out} get buried
under newer ones); transferring back from \cc{out} to \cc{in} after every dequeue
($\Theta(n)$ per operation).

\solution{Problem}{44}
\idea After enqueueing $x$, rotate the $n-1$ older items behind it (dequeue and
re-enqueue each). The front of the queue is then always the most recent item, the top.
\listingfile{sol/p44.c}{A stack from one queue.}
\complexity \cc{qsPush}: $\Theta(n)$; \cc{qsPop}, \cc{qsTop}: $\Theta(1)$.
\mistakes Evaluating \cc{qSize} after the enqueue and rotating all $n$ items, which puts
$x$ back at the rear; rotating in \cc{qsTop} as well as in \cc{qsPush}.

\solution{Problem}{45}
The solution is Listing~\ref{lst:josephus} of Chapter~\ref{ch:queues}.
\idea A queue is a circle: rotating $k-1$ people from the front to the rear and removing
the next one is one round of counting. The bonus uses the recurrence
$J(i)=(J(i-1)+k)\bmod i$ proved there.
\complexity Simulation $\Theta(nk)$; recurrence $\Theta(n)$ time, $O(1)$ space.
\mistakes Rotating $k$ people instead of $k-1$; returning the 0-based position of the
recurrence without adding 1; numbering from 0 in the simulation.

\solution{Problem}{46}
\idea A stack reverses: move the first $k$ items from the queue to the stack, then back
to the rear of the queue. They are reversed but now stand behind the other $n-k$ items,
so rotate those $n-k$ items once around.
\listingfile{sol/p46.c}{Reversing the first $k$ items of a queue.}
\complexity $\Theta(n)$: $k$ items move twice, $n-k$ once.
\mistakes Using \cc{qSize(q)} as the loop bound while the loop changes the queue (the
solution reads it into \cc{rest} first); rotating $n$ items instead of $n-k$.

\solution{Problem}{47}
\idea Pings expire in the order they arrived, so the recent ones form a queue: enqueue
\cc{t}, then dequeue from the front while the front is older than $t-3000$. The size of
the queue is the answer.
\listingfile{sol/p47.c}{Number of recent calls.}
\complexity $\Theta(1)$ amortised: each ping is enqueued once and dequeued at most once.
\mistakes Using \cc{<=} (the window is inclusive: a ping at exactly $t-3000$ counts);
scanning every stored ping on each call.

\solution{Problem}{48}
\idea With one stack and the queue itself, five passes. (1)~Push the first half;
(2)~pour the stack back to the rear (the first half is now reversed at the rear);
(3)~rotate the second half behind it; (4)~push the first half again --- the double
reversal restores its order on the stack, with $x_1$ on top; (5)~alternately pop one item
from the stack and rotate one item of the queue to the rear. The comments in the listing
show the queue for $[1..8]$ after each pass.
\listingfile{sol/p48.c}{Interleaving the halves of a queue.}
\complexity $\Theta(n)$: a constant number of passes over $n/2$ items.
\mistakes Pushing the first half once only, which interleaves it in reverse
(\cc{4 5 3 6 2 7 1 8}); using a second queue, which is not permitted.

\solution{Problem}{49}
\idea Keep a queue of the \emph{indices} of the negative numbers in the current window, in
order. When the window \cc{a[i-k+1..i]} is complete, drop from the front the indices that
have slid out; the front, if any, is the first negative of the window.
\listingfile{sol/p49.c}{First negative number in every window.}
\complexity $\Theta(n)$: every index is enqueued and dequeued at most once. The window
scan is $\Theta(nk)$.
\mistakes Storing values instead of indices (then it cannot be decided when an item
leaves the window); an off-by-one in the window start (\cc{front < i - k + 1}).
